\section{本章小结}
\begin{conclusion}
现在我们可以把一个三维点送到图像里，也可以用已知点与像素反求位姿。首先统一单位和坐标方向，再区分K/D、刚体变换与透视除法；求解返回成功后，还要看点序、深度和重投影，不让一个布尔值替你判断数据。

接下来是“怎样改得更好”：像素残差决定目标，参数化和扰动方向决定雅可比，增量方程提出候选，再用实际非线性结果判断。你应能运行固定点位姿优化与联合BA，说明后者为什么还要固定参考和尺度。低像素残差与高三维精度之间，并没有自动划等号。

RQ、SVD与指数映射各有自己的用途。三维对应已知时可以直接刚体拟合；对应未知时，ICP才反复找最近邻。公式推导不必第一次就全部背下，但程序里的方向、单位、误差与约束要能亲手核对，之后查文档换算法才有依据。
\end{conclusion}

\begin{problemset}
\item\label{prob:v2-ba-projection} 基础：fx=600、cx=300，X=0.1m、Z=2m，求u。整张图按坐标缩放0.5后，再从左侧裁10pixel，求新fx/cx及u。
\item\label{prob:v2-ba-frame} 基础：已知$q=Rp+t$，写出逆变换及相机光心在物体系中的坐标。若原始图已去畸变，为什么不能继续传原D再补偿一次？
\item\label{prob:v2-ba-pnp} 基础：说明校准DLT的像素预处理与非退化条件；三个P3P点为什么可能给多个位姿？第四点及正深度分别排除什么问题？
\item\label{prob:v2-ba-residual} 基础：观测(10,20)、预测(13,24)，按本章约定求残差及平方和。20点固定点优化与三相机联合项目的J尺寸分别是什么？写出正规方程右端的符号。
\item\label{prob:v2-ba-gauge} 提高：相机与世界点同时变换，为什么像素可以不变？解释固定第一相机后仍需尺度锚；Schur消元回代与丢掉旧变量各保留了什么？
\item\label{prob:v2-ba-matrix} 基础：M=KR，K上三角、R正交，应选哪种分解？若旋转投影得到行列式-1，怎样修正？说明修正后要核对的两个旋转条件。
\item\label{prob:v2-ba-graph} 提高：把BA路标固定，是否就得到相对位姿图？写出本章位姿图边的测量方向，说明为何协方差的逆才是信息矩阵。
\item\label{prob:v2-ba-icp} 基础：已知对应配准与ICP在哪一步不同？给一个“最近邻误差为零，但标签对应不对”的点云例子，说明还需要什么约束。
\item\label{prob:v2-ba-lie} 提高：旋转向量$(0,0,\pi)$作用于$(1,0,0)$的结果是什么？写出$[\rho,\phi]$顺序的帽矩阵，说明t何时不能直接当rho。
\item\label{prob:v2-ba-experiment} 实践：在同一生成条件下保存联合BA默认、无噪声、异常点普通平方与Huber四次结果。比较像素和三维误差，指出哪些结论只能来自无噪声对照、哪些不能由代价下降推出；用实际NPZ画同范围残差图。
\end{problemset}

\section{本章验收}
\begin{enumerate}
\item 检查者改变一个三维点、图像缩放或裁剪，你现场手算并运行basics核对；展示变换链的方向、逆变换与单位，再演示PnP点序/尺度变体及全部P3P候选检查。
\item 展示RQ重建、旋转约束与SO(3)/SE(3)指数检查，运行独立差分核对像素雅可比，指出残差符号、变量顺序及左右扰动如何保持一致。
\item 运行固定点位姿优化与三相机联合BA，保留输入、实际数组及误差图；展示异常点对照，解释参考/尺度锚、71变量/120残差和Schur回代误差，用实际误差说明结果边界。
\item 演示已知对应拟合、真实最近邻ICP与一维闭环变体；说明各自的边或对应是什么、何时不唯一，以及停止条件为何不能替代物理对应检查。
\end{enumerate}
本章验收使用声明条件的软件几何实验，实际车辆坐标安装、时序及赛场精度仍按卷三任务核验。不能拿这里的合成数值当作相机实测精度，也不能只交一张“优化成功”的截图。
